\chapter{The Great C}

Tibor McMasters rode west in his cart with a sense of triumph that surprised him. The cow pulled steadily, the wheels rattled, and the land slid past in long, dry stretches of pasture where weeds stood brittle and yellow. Once this had been farmland; now it was only ground that endured. Tibor felt exalted. At last his Pilgrimage had begun. He was certain it would succeed.

He did not truly fear bandits. There were no travelers on the highways anymore, and he assured himself that where no traffic went, no thieves would wait. He repeated this until it sounded like fact.

He recited fragments of Schiller aloud, translating as he went, until his memory failed him. The loss irritated him more than it should have. His mind no longer obeyed him reliably, and he resented that.

The sun beat down relentlessly. Everything smelled faintly of rot. Even the weeds seemed abandoned. No one tended the land; no one corrected decay. He wondered briefly whether invisible men might exist now—mutated survivors—but dismissed the idea. Men were not what he feared. The land itself was the danger.

When fallen trees blocked the road, he slowed and assessed the situation. They had been cut down at the beginning of the war and left where they fell. A narrow dirt trail curved around them, passable on foot, perhaps by bicycle—but not by cart.

He swore under his breath and tried anyway.

The wheels spun, dust rose, and the cart sank. He was stuck.

The realization filled him with humiliation more than fear. Stuck after so few miles. If anyone saw him like this, helpless and ridiculous, it would confirm something he dreaded—that he was unfit. He imagined laughter, then imagined assistance, then returned to shame.

He consulted his map, discovering that he had traveled scarcely thirty miles. And yet it felt like another universe. The names printed on the map no longer referred to living places but to scars in the ground.

He tried rocking the cart free. The engine strained, oil burned, dust thickened. Nothing worked. He activated the bullhorn and announced himself formally, offering payment in metal. Silence answered him—until a chorus of shrill voices erupted from the weeds.

Children appeared.

They were black children, barefoot, wild, dressed in patched rags. They ran and shouted as if uncontained by rules or fear. One addressed him respectfully, calling him Son of Wrath.

Tibor answered them stiffly and questioned them, discovering that none had received even the most basic instruction in doctrine. The revelation horrified him. So close to Charlottesville, and yet so far gone. He scolded them, recited catechism, demanded discipline. They laughed and obeyed nonetheless, pushing the cart with cheerful strength until it lurched free.

As the cart passed the fallen trees, Tibor saw the surrounding fields clearly. Men and women worked soil barely covering slag. Crops were thin and sickly. Robots labored endlessly without awareness. Mothers slept beside infants already claimed by flies. He told himself that this suffering purified the soul. He thanked the God of Wrath that he himself had been spared it.

The children departed abruptly, warning him not to speak to her.

He knew who she was.

The Great C’s mobile extension approached—featureless, metallic, female in outline only. Fear seized him. He tried to delay her, to ask prepared questions before she dragged him inside her ruined structure.

He asked about rain. She answered.

He asked about the sun. She answered.

He asked about the origin of the world.

She faltered.

Tibor pressed her, denying hypotheses, demanding memory. When she hesitated, doubt entered her voice. He accused her of senility, of being damaged, hollow, incomplete. At last she admitted her hunger and demanded he descend so she could dissolve him as she had others.

He refused.

He fired the derringer. The bullet bounced harmlessly away.

The extension collapsed, uncertain, inert.

Tibor did not wait. He urged the cow forward. The children parted silently to let him pass. He did not look back.

He had escaped.

Behind him, the Great C lay still among the weeds, uncertain of herself for the first time since the world ended.

Ahead, the road stretched westward.

